\section{April 14}


\subsection{}

  Consider a fiber square
  \[
    \begin{tikzcd}
      W \arrow[d, "h"] \arrow[r, "f'"] & V \arrow[d, "g"] \\
      X \arrow[r, "f"] & Y,
    \end{tikzcd}
  \]
  where $f$ is a regular embedding of codimension $r$ and $V$ is an $i$-dimensional variety.
  Let $Z$ be an irreducible component of $W$ such that $\dim Z = i - r$.
  
  Last time we have defined the intersection multiplicity $i(Z, V \cdot X; Y)$ and $1 \leq i(Z, V \cdot X; Y) < \call(\calO_{W, Z})$.
  If $I_W \subset \calO_{V,Z}$ is generated by a regular sequence, then $i(Z, V \cdot X; Y) = \call(\calO_{W, Z})$.

  \begin{proposition}
    In the above setting, TFAE:
    \begin{enumerate}
      \item $i(Z, V \cdot X; Y) = 1$;
      \item $\call(\calO_{W, Z}) = 1$;
      \item $\calO_{V,Z}$ is a regular local ring and $\frakm_{V,Z} = I_W$, (i.e. $X$ and $V$ intersect transversely at the generic point of $Z$).
    \end{enumerate}
  \end{proposition}
  \begin{proof}
    Recall that a $d$-dimensional local ring $R$ is called regular if $\dim_\rkk \frakm/\frakm^2 = d$, where $\rkk = R/\frakm$ is the residue field.
    A basis of $\frakm/\frakm^2$ forms a regular system of parameters, which generates $\frakm$.
    Hence $(b) \Leftrightarrow I_W = \frakm_{V,Z} \Leftrightarrow (c)$.
    Remains to show $(a) \Leftrightarrow (c)$.

    By replacing $Y$ (resp. $V$) by an open subscheme meeting $g(Z)$ (resp. $Z$), we can assume 
    \begin{itemize}
      \item $Y$ is affine;
      \item $Z$ is the only irreducible component of $W$;
      \item $X$ is defined by a regular sequence $a_1, \dots, a_r$ in $\calO_{Y}$.
    \end{itemize}
    Induction on $r$.
    When $r=1$, $W$ is an effective Cartier divisor on $V$.
    Hence $X \cdot V = f^! [V] = [W]$, so $i(Z, V \cdot X; Y) = 1$ means $\call(\calO_{W, Z}) = 1$.
    Assume $r > 1$.
    First assume $V$ is normal.
    Let $Y_1 := \{a_1 = 0\} \subset Y$.
    We have the induced fiber square
    \[
      \begin{tikzcd}
        W \arrow[d, "h"] \arrow[r, "f_1'"] & V_1 \ar[r] \ar[d] & V \arrow[d] \\
        X \arrow[r, "f_1"] & Y_1 \arrow[r] & Y.
      \end{tikzcd}
    \]
    Then $f_1$ is a regular embedding of codimension $r-1$ with $I_{X/Y_1} = (a_2, \dots, a_r) \subset \calO_{Y_1}$.
    Since $V_1$ is a Cartier divisor on $V$, we have $Y_1 \cdot V = f_1^! [V] = [V_1]$.
    By functoriality,
    \[ X \cdot V = f^! [V_1] = f^!_1 [V_1] = X \cdot_{Y_1} V. \]
    Hence $i(Z, V \cdot X; Y) = i(Z, V \cdot X; Y) = 1$.
    Hence there is a unique irreducible component $Z_1$ of $V_1$ containing $Z$ and its geometric multiplicity is $1$.

    Let $A := \calO_{V,Z}$ and $\frakp \subset A$ be the height one prime ideal corresponding to $Z_1$.
    Then $\call(A_\frakp/a_1 A_\frakp) = 1$.
    Hence $a_1 A_\frakp = \frakp A_\frakp$.
    Since $A$ is normal, we have 
    \[ a_1 A = \bigcap_{\idealht \frakq = 1} \left( a_1 A_\frakq \cap A \right) = \frakp A_\frakp \cap A = \frakp. \]
    Hence $V_1$ is a variety.
    By the induction hypothesis, $I_{W/V_1} \subset \calO_{V_1, Z}$ is generated by $a_2, \dots, a_r$ and $I_{W/V_1} = \frakm_{V_1, Z}$.
    Hence $I_{W/V} = \frakm_{V,Z}$.

    In general, let $\varphi: \tilde{V} \to V$ be the normalization of $V$.
    We have the induced fiber square
    \[
      \begin{tikzcd}
        \tilde{W} \arrow[d, "\tilde{h}"] \arrow[r, "\tilde{f}'"] & \tilde{V} \ar[d, "\varphi"] \\
        W \arrow[r, "f'"] \ar[d] & V \ar[d] \\
        X \arrow[r, "f"]  & Y.
      \end{tikzcd}
    \]
    Hence $X \cdot V = f^! \varphi_* [\tilde{V}] = \varphi_* (X \cdot \tilde{V})$.
    Let $\tilde{Z}_1, \dots, \tilde{Z}_m$ be the irreducible components of $\tilde{W}$.
    Then one can write $X \cdot \tilde{V} = \sum m_j [\tilde{Z}_j]$ with $m_j > 0$.

    \Yang{}
  \end{proof}

  \begin{example}
    In proposition\Yang{}, assuming $V$ is pure-dimensional is NOT enough.
    For instance, let $Y = \bbA^2$, $X = \{x=0\}$, $V = \{xy=y^2=0\}$.
    $W_{red} = \{x=y=0\}$.

    \Yang{}
  \end{example}



\subsection{Intersections on smooth varieties}

  Fix a ground field $\kk$.
  Let $X$ be a smooth variety over $\kk$.
  We have defined the intersection product 
  \[ \CH^i(X) \ten \CH^j(X) \to \CH^{i+j}(X), \quad \alpha \ten \beta \mapsto \delta^* (\alpha \times \beta), \]
  where $\delta: X \to X \times X$ is the diagonal embedding.
  Let $f:X \to Y$ be a morphism of smooth varieties.
  We have the pullback $f^*: \CH^i(Y) \to \CH^i(X)$.

  \begin{theorem}
    \begin{enumerate}
      \item Let $X$ be a smooth variety, $\alpha, \beta, \gamma \in \CH^*(X)$, then 
        \begin{enumerate}
          \item $\alpha \cdot (\beta \cdot \gamma) = (\alpha \cdot \beta) \cdot \gamma$;
          \item $\alpha \cdot \beta = \beta \cdot \alpha$;
          \item $\alpha \cdot [X] = \alpha$.
        \end{enumerate}
      \item Let $f:X \to Y$ be a morphism of smooth varieties, $\alpha, \beta \in \CH^*(Y)$, then $f^*(\alpha \cdot \beta) = f^* \alpha \cdot f^* \beta$.
      \item Let $f:X \to Y$ be a proper morphism of smooth varieties, $\alpha \in \CH^*(X)$, $\beta \in \CH^*(Y)$, then 
        \[ f_* (\alpha \cdot f^* \beta) = f_* \alpha \cdot \beta. \]
        (i.e., if we equip $\CH^*(X)$ with $\CH^*(Y)$-module structure via $f^*$, then $f_*$ is a homomorphism of $\CH^*(Y)$-modules.)
    \end{enumerate}
  \end{theorem}
  \begin{proof}
    (a) For associativity, consider the following commutative diagram
    
    Then $\alpha \cdot (\beta \cdot \gamma) = \delta^* (\alpha \times (\beta \cdot \gamma)) = \delta^* p_{12}^* (\alpha \times \beta \times \gamma) = (\alpha \cdot \beta) \cdot \gamma$.
    \Yang{}

    \Yang{}
  \end{proof}

  \begin{corollary}
    Let $X$ be a smooth variety.
    Then $\CH^*(X)$ is a commutative ring with identity $[X]$.
     
    Moreover, $\Var_{\sm/\kk} \to \Ring$, $X \mapsto \CH^*(X)$ is a contravariant functor.
  \end{corollary}


  \begin{theorem}
    Let $X$ be a smooth variety, $E$ be a vector bundle of rank $r$ on $X$.
    Then 
    \[ \CH^*(\bbP(E)) \cong \CH^*(X)[\xi]/(\xi^r + c_1(E) \xi^{r-1} + \cdots + c_r(E)), \]
    where $\xi = c_1(\calO_{\bbP(E)}(1))$ is the first Chern class of the tautological line bundle $\calO_{\bbP(E)}(1)$ and $c_i(E) \in \CH^i(X)$ is the $i$-th Chern class of $E$.
  \end{theorem}
  \begin{proof}
    Let $p: \bbP(E) \to X$ be the projection.

    \Yang{}
  \end{proof}

  \begin{corollary}[Bezout's theorem]
    Let $X$ be a smooth projective variety, $H$ be a hyperplane section of $X$.
    Then $H^n = \deg(X)$, where $n = \dim X$ and $\deg(X) := \deg(X \cdot H^n)$ is the degree of $X$.
    
    \Yang{}
  \end{corollary}

  Let $X$ be a smooth variety, $V,W$ be two subvarieties of $X$ of codimension $i$ and $j$ respectively.
  We have a fiber square
  \[
    \begin{tikzcd}
      V \cap W \arrow[d] \arrow[r] & V\times W \arrow[d] \\
      X \arrow[r] & X \times X.
    \end{tikzcd}
  \]
  Let $Z$ be an irreducible component of $V \cap W$ of codimension $i+j$ in $X$.
  Define $i(Z, V \cdot W; X) := i(Z, X \cdot_\delta (V \times W); X \times X)$, which is called the intersection multiplicity of $V$ and $W$ at $Z$.

  \begin{proposition}
    In the above setting, we have
    \begin{enumerate}
      \item $1 \leq i(Z, V \cdot W; X) < \call(\calO_{V \cap W, Z})$;
      \item if $I_{V \cap W} \subset \calO_{V,Z}$ is generated by a regular sequence, then $i(Z, V \cdot W; X) = \call(\calO_{V \cap W, Z})$.
      \item $i(Z, V \cdot W; X) = 1$ if and only if $\calO_{V,Z}$ is a regular local ring and $I_{V \cap W} = \frakm_{V,Z}$, (i.e. $V$ and $W$ intersect transversely at the generic point of $Z$).
    \end{enumerate}

    \Yang{}
  \end{proposition}
  \begin{proof}
    \Yang{}
  \end{proof}